% Irodov Part IV §4.1 independent original problems 66--81.
% User PDF pp.172--175, printed pp.163--165, original answers PDF 325--326.
\oscproblem{66}{163}
两根竖直铜轨之间距为 $\ell$，
上端由自感为 $L$ 的线圈短接，
水平活动导杆质量为 $m$，
起初静止，随后沿两导轨自由下落。
整个装置处于垂直导轨平面的匀强磁场 $B$ 中，
各导体电阻忽略。
求导杆下降位移 $x(t)$。
\exsolution
取向下为 $x$ 正向，初始 $x(0)=0$、
$\dot x(0)=0$，初始线圈电流为零。
磁链守恒
$LI+B\ell x=0$，
所以 $I=-(B\ell/L)x$。
活动杆的安培力沿上方、大小
$B^2\ell^2x/L$。
运动方程
\[
m\ddot x=mg-\frac{B^2\ell^2}{L}x.
\]
令
$\omega_0=B\ell/\sqrt{mL}$，
方程化为
$\ddot x+\omega_0^2x=g$。
因此
\[
\boxed{x(t)=\frac{g}{\omega_0^2}
(1-\cos\omega_0t)}.
\]
\excheck 从静止起初加速度为 $g$，
因 $I(0)=0$；回路无电阻时系统能量只在重力势、动能与磁能间交换。

\oscproblem{67}{163}
质点的阻尼振动规律为
$x=a_0e^{-\beta t}\sin\omega t$，
其中 $\beta,\omega>0$。
求 (a) $t=0$ 时位移振幅和速度；
(b) 各个位置转折点所对应的时刻。
\exsolution
幅度包络为
$A(t)=a_0e^{-\beta t}$，
初始包络为
$\boxed{A(0)=a_0}$。
速度
$v=a_0e^{-\beta t}[\omega\cos\omega t-
\beta\sin\omega t]$，
故
$\boxed{v(0)=a_0\omega}$。
转折点满足 $v(t_n)=0$，
即
$\tan\omega t_n=\omega/\beta$，
因此
\[
\boxed{t_n=\frac1\omega
\left[\arctan\frac{\omega}{\beta}+n\pi\right],
\quad n=0,1,2,\ldots}.
\]
\excheck 阻尼会使转折点提前于普通正弦振动的
$\omega t=\pi/2$；$\beta\to0$ 时恢复其极限。

\oscproblem{68}{163}
刚体作阻尼扭转振动，
$\varphi(t)=\varphi_0e^{-\beta t}\cos\omega t$。
求 (a) 初始角速度和角加速度；
(b) 角速度达到局部极值的时刻。
\exsolution
直接求导：
\[
\dot\varphi=-\varphi_0e^{-\beta t}
(\beta\cos\omega t+\omega\sin\omega t),
\]
\[
\ddot\varphi=\varphi_0e^{-\beta t}
[(\beta^2-\omega^2)\cos\omega t+
2\beta\omega\sin\omega t].
\]
所以
\[
\boxed{\dot\varphi(0)=-\beta\varphi_0},
\quad
\boxed{\ddot\varphi(0)=(\beta^2-\omega^2)\varphi_0}.
\]
角速度极值条件 $\ddot\varphi(t)=0$，
\[
\boxed{t_n=\frac1\omega
\left[\arctan\frac{\omega^2-\beta^2}{2\beta\omega}
+n\pi\right]},
\]
其中选取使 $t_n\ge0$ 的整数 $n$。
\excheck 这些是\emph{有符号角速度}的局部极值点；
比较速度模的全部峰值时需区分正峰与负峰。

\oscproblem{69}{163}
阻尼振动写成
$x(t)=a_0e^{-\beta t}\cos(\omega t+\alpha)$。
求初始振幅 $a_0$、初相 $\alpha$，
当初始条件为
(a) $x(0)=0$、$\dot x(0)=v_0$；
(b) $x(0)=x_0$、$\dot x(0)=0$。
\exsolution
一般初值满足
$x_0=a_0\cos\alpha$、
$v_0=-a_0(\beta\cos\alpha+\omega\sin\alpha)$。
(a) $\cos\alpha=0$，
$a_0=|v_0|/\omega$，
并由 $v_0=-a_0\omega\sin\alpha$ 决定相位，
\[
\boxed{a_0=\frac{|v_0|}{\omega},
\quad
\alpha=
\begin{cases}
-\pi/2,&v_0>0,\\
+\pi/2,&v_0<0.
\end{cases}}
\]
(b) $v_0=0$ 给
$\tan\alpha=-\beta/\omega$；
由 $x_0=a_0\cos\alpha$，
\[
\boxed{a_0=|x_0|
\sqrt{1+\left(\frac\beta\omega\right)^2}}.
\]
当 $x_0>0$，
可取
$\alpha=-\arctan(\beta/\omega)$；
当 $x_0<0$，
按相位加 $\pi$ 选择等效分支。
\excheck $\beta=0$ 时退化为标准无阻尼简谐振动初值关系。

\oscproblem{70}{163}
某阻尼振动角频率
$\omega=25\,\mathrm{s^{-1}}$，
在 $t=0$ 时瞬时速度为零，
而实际位移模比该时刻的包络振幅小
$\eta=1.020$ 倍。
求阻尼系数 $\beta$。
\exsolution
写
$x=a_0e^{-\beta t}\cos(\omega t+\alpha)$。
初始速度零给
$\beta\cos\alpha+\omega\sin\alpha=0$，
即
$\tan\alpha=-\beta/\omega$。
位移与包络幅度之比
$|x(0)|/a_0=|\cos\alpha|
=1/\sqrt{1+(\beta/\omega)^2}$。
题设此比为 $1/\eta$，所以
\[
\boxed{\beta=\omega\sqrt{\eta^2-1}
=25\sqrt{1.020^2-1}
\simeq5.0\,\mathrm{s^{-1}}}.
\]
\excheck 对阻尼振动，“速度为零的转折点”
并不恰等于瞬时包络振幅，
其差由 $\beta/\omega$ 决定。

\oscproblem{71}{164}
质点阻尼振动角频率为 $\omega$、
阻尼系数为 $\beta$。
求速度振幅包络随时间的关系，
当 $t=0$：
(a) 位移幅度为 $a_0$；
(b) 实际位移为零，实际速度为 $v_0$。
\exsolution
一般位移
$x=a_0e^{-\beta t}\cos(\omega t+\alpha)$，
速度可写成
\[
v=-a_0e^{-\beta t}\sqrt{\omega^2+\beta^2}
\cos(\omega t+\alpha-\delta)
\]
（适当取相位 $\delta$）。
故
\[
\boxed{V_a(t)=a_0\sqrt{\omega^2+\beta^2}\,
e^{-\beta t}}.
\]
(b) 初始 $x(0)=0$，
按第69题有
$a_0=|v_0|/\omega$，
从而
\[
\boxed{V_b(t)=|v_0|
\sqrt{1+\left(\frac\beta\omega\right)^2}
e^{-\beta t}}.
\]
\excheck 速度包络与位移包络的比例是
$\sqrt{\omega^2+\beta^2}$，不是简单的 $\omega$。

\oscproblem{72}{164}
两种阻尼振动的参数为
$T_1=0.10\,\mathrm{ms}$、
$\beta_1=100\,\mathrm{s^{-1}}$；
$T_2=10\,\mathrm{ms}$、
$\beta_2=10\,\mathrm{s^{-1}}$。
判断哪一种衰减更快。
\exsolution
振幅包络 $A(t)=A_0e^{-\beta t}$。
若比较\emph{相同真实时间}，
$\beta_1>\beta_2$，
所以第一种衰减更快，
其时间常数分别为
$\tau_1=1/\beta_1=0.010$s、
$\tau_2=1/\beta_2=0.10$s。
若比较\emph{每振动一周}衰减程度，
对数减量
$\lambda_i=\beta_iT_i$：
\[
\boxed{\lambda_1=100(10^{-4})=0.010},
\qquad
\boxed{\lambda_2=10(10^{-2})=0.10}.
\]
因此按单周期比例衰减，第二种更快。
\excheck 题目未给“更快”的比较基准，
两个答案分别对应真实时间和自身周期，
与原书书末的提醒一致。

\oscproblem{73}{164}
单摆在某介质中振荡，
对数减量为
$\lambda_0=1.50$。
若介质阻力系数加大
$n=2.00$ 倍，求新对数减量；
阻力加大多少倍后不再出现振荡？
\exsolution
设无阻尼固有角频率为
$\omega_0$，
原阻尼系数 $\beta_0$，
阻尼振动频率
$\omega_{d0}=\sqrt{\omega_0^2-\beta_0^2}$。
对数减量
$\lambda_0=2\pi\beta_0/\omega_{d0}$。
阻力倍数为 $n$ 时
$\beta=n\beta_0$，
$\omega_d=\sqrt{\omega_0^2-n^2\beta_0^2}$，
故
\[
\boxed{\lambda(n)=
\frac{n\lambda_0}
{\sqrt{1+(1-n^2)(\lambda_0/2\pi)^2}}}.
\]
代 $n=2$ 得
$\boxed{\lambda\simeq3.3}$。
临界阻尼为
$n_c\beta_0=\omega_0$，
因此
\[
\boxed{n_c=\sqrt{1+\left(\frac{2\pi}{\lambda_0}\right)^2}
\simeq4.3}.
\]
\excheck $n\ge n_c$ 时系统转为临界或过阻尼，
不再围绕平衡位置反复振荡。

\oscproblem{74}{164}
质量块悬挂在无质量轻弹簧下，
静伸长
$\Delta x=9.8$cm；
轻微竖直推动后阻尼振荡的对数减量为
$\lambda=3.1$。
求实际振荡周期。
\exsolution
静平衡
$k\Delta x=mg$，
无阻尼固有角频率
$\omega_0^2=k/m=g/\Delta x$。
对于欠阻尼振荡，
$\omega_d^2=\omega_0^2-\beta^2$，
周期 $T=2\pi/\omega_d$，
对数减量
$\lambda=\beta T$。
消去 $\beta$：
\[
\frac{4\pi^2}{T^2}
=\frac g{\Delta x}-\frac{\lambda^2}{T^2},
\]
故
\[
\boxed{T=
\sqrt{\frac{(4\pi^2+\lambda^2)\Delta x}{g}}
\simeq0.70\,\mathrm s}.
\]
\excheck 若错误采用无阻尼周期
$2\pi\sqrt{\Delta x/g}$，
将忽略 $\lambda=3.1$ 的显著阻尼效应。

\oscproblem{75}{164}
某振子经过
$n=110$ 次振动后，位移幅度降低
$\eta=2.0$ 倍。
求品质因数 $Q$。
\exsolution
若每周期振幅比
$A_{k+1}/A_k=e^{-\lambda}$，
则经过 $n$ 周期有
$A_n/A_0=e^{-n\lambda}=1/\eta$，
所以
$\lambda=(\ln\eta)/n$。
欠阻尼振子的品质因数按
$Q\simeq\pi/\lambda$，故
\[
\boxed{Q\simeq\frac{\pi n}{\ln\eta}
=\frac{\pi(110)}{\ln2}
\simeq5.0\times10^2}.
\]
\excheck 这里 $\lambda\ll1$，
因此用 $Q\simeq\pi/\lambda$ 的弱阻尼近似非常好。

\oscproblem{76}{164}
粒子初始离开平衡位置 $\ell=1.0$cm，
随后由静止释放，
发生阻尼振动直至停止。
对数减量
$\lambda=0.020$。
求整个过程中累计走过的总路程。
\exsolution
相邻两个转折点间隔半个阻尼周期，
对应位移极值的模按因子
$r=e^{-\lambda/2}$ 递减。
起始极值为 $\ell$，
之后转折点依次为
$-\ell r,+\ell r^2,\ldots$。
总路径为
\[
s=\ell+\ell r+
(\ell r+\ell r^2)+(\ell r^2+\ell r^3)+\cdots
=\ell\frac{1+r}{1-r}.
\]
因此
\[
\boxed{s=\ell
\frac{1+e^{-\lambda/2}}{1-e^{-\lambda/2}}
\simeq2.0\,\mathrm m}.
\]
\excheck 虽然每次反转位移很小，
衰减率仅每半周期约1%，
累计路程可远大于初始1cm振幅。

\oscproblem{77}{164}
数学单摆长度 $\ell=50$cm，
在
$\tau=5.2$min 内总机械能下降
$\eta=4.0\times10^4$ 倍。
求品质因数。
\exsolution
阻尼振动能量包络
$W(t)\propto e^{-2\beta t}$，
故
$\eta=e^{2\beta\tau}$，
$\beta=(\ln\eta)/(2\tau)$。
无阻尼固有角频率
$\omega_0=\sqrt{g/\ell}$，
阻尼角频率
$\omega_d=\sqrt{\omega_0^2-\beta^2}$。
用本书弱阻尼定义
$Q=\omega_d/(2\beta)$，
\[
\boxed{Q=\frac12
\sqrt{\frac{4g\tau^2}{\ell(\ln\eta)^2}-1}
\simeq1.3\times10^2}.
\]
\excheck $5.2$min 需换为312s；
$Q\gg1$ 支持弱阻尼模型。

\oscproblem{78}{164}
均匀薄圆盘半径
$R=13$cm，
可绕垂直盘面、穿过盘缘的水平轴转动，
作小幅阻尼摆动，对数减量
$\lambda=1.00$。
求振荡周期。
\exsolution
圆盘对悬轴的转动惯量
$J=J_{\rm cm}+mR^2=(3/2)mR^2$。
质心距悬轴 $R$，
故无阻尼角频率
$\omega_0^2=mgR/J=2g/(3R)$。
对数减量 $\lambda=\beta T$，
阻尼角频率
$\omega_d=2\pi/T$，
于是
$\omega_0^2=(4\pi^2+\lambda^2)/T^2$，
\[
\boxed{T=
\sqrt{\frac{3R(4\pi^2+\lambda^2)}{2g}}
\simeq0.90\,\mathrm s}.
\]
\excheck 悬轴在\emph{盘缘}且轴垂直盘面，
所以 $J$ 不是圆盘中心惯量 $mR^2/2$。

\oscproblem{79}{164}
水平薄均匀圆盘质量为 $m$、半径为 $R$，
由弹性细线悬挂，扭转回复力矩大小
$\alpha|\varphi|$。
浸于黏性流体中，
每单位盘面面积的阻力大小
$\eta v$（$v$ 为局部速度）。
求小扭转振动角频率。
\exsolution
圆盘绕轴转动惯量
$J=mR^2/2$。
半径 $r$ 处的速度
$v=r|\dot\varphi|$，
每个表面元受到切向阻力密度
$\eta r|\dot\varphi|$，
相对轴的力矩元
$\eta r^2\dot\varphi\,\mathrm dS$。
盘上下两面总阻力矩系数
\[
b=2\eta\int_{\rm disc}r^2\,\mathrm dS
=2\eta\left(\frac{\pi R^4}{2}\right)
=\pi\eta R^4.
\]
方程
$J\ddot\varphi+b\dot\varphi+\alpha\varphi=0$。
设阻尼系数
$\beta=b/(2J)=\pi\eta R^2/m$，
则
\[
\boxed{\omega=
\sqrt{\frac{2\alpha}{mR^2}-
\left(\frac{\pi\eta R^2}{m}\right)^2}}.
\]
\excheck 只有根号内为正时存在欠阻尼振动；
两面流体黏性阻力均需计入。

\oscproblem{80}{165}
图4.24：半径为 $R$ 的薄盘 A
在上下两个固定平面之间扭转振动，
盘对中心轴转动惯量为 $J$，
盘与每个平面的间隙为
$h\ll R$。
振动周期 $T$、对数减量 $\lambda$ 已知。
求周围气体的动力黏度 $\eta$。
\exsolution
半径 $r$ 处盘面切向速度
$v=r\dot\varphi$，
在狭缝内近似线性速度剖面，
切应力为
$\tau_s=\eta v/h=\eta r\dot\varphi/h$。
上下两面总阻力矩
\[
M_{\rm drag}=-2\frac{\eta\dot\varphi}{h}
\int_{\rm disc}r^2\,\mathrm dS
=-\frac{\pi\eta R^4}{h}\dot\varphi.
\]
即线性阻尼系数
$b=\pi\eta R^4/h$。
运动方程
$J\ddot\varphi+b\dot\varphi+k\varphi=0$，
位移包络衰减率
$\beta=b/(2J)$，
对数减量 $\lambda=\beta T$。
所以
\[
\boxed{\eta=\frac{2\lambda hJ}{\pi R^4T}}.
\]
\excheck $\eta$ 维度为
$\mathrm{Pa\,s}$；
本式要求层流薄隙 $h\ll R$ 以及线性黏滞近似。

\oscproblem{81}{165}
图4.25：边长为 $a$ 的正方形导线框
悬在弹性细线下，
处于均匀水平磁场 $B$ 中。
平衡时框面平行磁场，
绕通过框中心的竖轴作微小扭转振动，
框对轴惯量为 $J$、电阻为 $R$，
忽略自感。
求振动角幅度衰减到原来的
$1/e$ 所需时间。
\exsolution
微小转角 $\varphi$ 时，
框面法向磁通
$\Phi=Ba^2\sin\varphi
\simeq Ba^2\varphi$。
感生电动势
$\mathcal E=-Ba^2\dot\varphi$，
感生电流
$I_{\rm ind}=\mathcal E/R$。
磁力矩
$\tau_m=I_{\rm ind}\,\mathrm d\Phi/\mathrm d\varphi
\simeq-(B^2a^4/R)\dot\varphi$，
等效黏性阻尼系数
$b=B^2a^4/R$。
微小扭转方程
$J\ddot\varphi+b\dot\varphi+k\varphi=0$，
在欠阻尼条件下振幅包络为
$A(t)=A_0e^{-\beta t}$，
$\beta=b/(2J)$。
所以
\[
\boxed{\tau_{1/e}=\frac1\beta
=\frac{2JR}{a^4B^2}}.
\]
\excheck 题设未给扭转弹性常数 $k$，
但幅度包络的 e 倍衰减时间只由
惯量和线性电磁阻尼系数决定。
